% ATTEMPT_STATUS: unresolved
% TOP_PROBLEM: {"problemId": "problem.hodge-conjecture", "canonicalTitle": "Hodge Conjecture", "releaseRank": 4, "primaryDomain": "geometry_topology", "primaryDomainLabel": "Geometry & topology", "displayStatus": "open", "releaseStatus": "open"}
% !TeX program = lualatex
% Research continuation by GPT 6 Astra Ultra, 27 September 2026.
% Root statement and inherited source list preserved; new work remains unresolved.
\documentclass[11pt]{article}
\usepackage[margin=25mm]{geometry}
\usepackage{fontspec}
\setmainfont{Latin Modern Roman}
\usepackage{amsmath,amssymb,mathtools}
\usepackage{unicode-math}
\setmathfont{Latin Modern Math}
\usepackage{xurl,hyperref}
\hypersetup{colorlinks=true,urlcolor=blue}
\setcounter{secnumdepth}{0}
\setlength{\parindent}{0pt}
\setlength{\parskip}{5pt}
\setlength{\emergencystretch}{3em}
\begin{document}
% problemId: problem.hodge-conjecture
% source releaseRank: 4; source releaseStatus: open
\section*{\texorpdfstring{Hodge Conjecture}{Hodge Conjecture}}\label{4}
\subsection{Definitions and mathematical statement}
Let $X$ be a smooth projective complex variety of dimension $d$. Singular cohomology has the Hodge decomposition $H^k(X,{\mathbb{C}})=\bigoplus_{a+b=k}H^{a,b}(X)$. Define
\[
 \operatorname{Hdg}^{p}(X)=H^{2p}(X,{\mathbb{Q}})\cap H^{p,p}(X),\qquad 0\le p\le d.
\]
A rational codimension-$p$ algebraic cycle is a finite sum $z=\sum_j a_j Z_j$, with $a_j\in{\mathbb{Q}}$ and $Z_j\subseteq X$ closed irreducible subvarieties of codimension $p$. Its class is $\operatorname{cl}(z)=\sum_j a_j[Z_j]$. The conjecture says
\[
 \operatorname{Hdg}^{p}(X)=\operatorname{span}_{{\mathbb{Q}}}\{[Z]:\operatorname{codim}_X Z=p\}
\]
for all $X,p$. Integral coefficients and nonprojective compact Kähler manifolds are not the selected formulation.
\par\smallskip\textit{Formulation and concept references:} \href{https://www.claymath.org/library/monographs/MPPc.pdf}{[S1]}, \href{https://www.claymath.org/wp-content/uploads/2022/06/hodge.pdf}{[E1]}.

\subsection{Short English statement}
Does every rational cohomology class of the appropriate Hodge type come from a rational combination of algebraic subvarieties?
\subsection{Research attempt}

\paragraph{Scope and literature checkpoint.}
This attempt concerns rational classes on smooth projective complex varieties.
Deligne's formulation \href{https://www.claymath.org/wp-content/uploads/2022/06/hodge.pdf}{[E1]}, Sections 1--2, supplies the divisor case and
explains why changing to integral coefficients changes the question. No claim
listed in the inherited catalogue is being used to settle the general case.
The work below tests two tempting constructions: generating cycles from
divisors, and replacing a Hodge class by an effective cycle. Both are stronger
than the requested conclusion, and both fail on an explicit class that is
itself algebraic. The point is to locate a precise obstruction before trying
to extend either construction. Historical novelty of these observations is
not claimed.

\paragraph{Attempt A: propagate the divisor case by hyperplane sections.}
Write $L(\alpha)=h\smile\alpha$ for an ample integral class on a connected
$d$-fold. If $d\ge2$, hard Lefschetz gives the rational isomorphism
\[
 L^{d-2}:H^2(X,\mathbb Q)\longrightarrow H^{2d-2}(X,\mathbb Q).
\]
The map shifts Hodge bidegrees by $(d-2,d-2)$. Thus for a rational
$(d-1,d-1)$ class $\gamma$, its unique inverse $\beta$ is rational and of
type $(1,1)$. The divisor theorem supplies a rational divisor class for
$\beta$; intersecting it with $d-2$ hyperplanes represents $\gamma$.
These standard inputs are described in \href{https://www.claymath.org/wp-content/uploads/2022/06/hodge.pdf}{[E1]}, Sections 2 and 4. Along with
$p=0,d$, this proves the conjecture in every codimension when $d\le3$.

This calculation suggests applying inverse Lefschetz operators until every
class is reduced to a divisor. The obstruction is already linear algebra.
For $2p\le d$, define
\[
 P^{2p}(X)=\ker\bigl(L^{d-2p+1}:H^{2p}(X,\mathbb Q)
                       \longrightarrow H^{2d-2p+2}(X,\mathbb Q)\bigr).
\]
The Lefschetz decomposition has a summand $P^{2p}(X)$ as well as the
summands obtained by multiplying lower-degree primitive classes by $h$.
An inverse on the image of $L$ does not remove this new summand. In
particular, on a fourfold, primitive degree-four Hodge classes cannot be
disposed of by the degree-two divisor theorem. Declaring an inverse
cohomological operation to be an algebraic correspondence would itself
require an additional algebraicity theorem. The subsequent argument does
not assume one.

\paragraph{Attempt B: test the entire divisor-generated algebra.}
Let $S$ be any smooth projective complex K3 surface and put $X=S\times S$.
We use the following standard K3 facts from \href{https://www.math.uni-bonn.de/people/huybrech/K3Global.pdf}{[E2]}, Chapters 1 and 3:
$H^1(S,\mathbb Q)=H^3(S,\mathbb Q)=0$, $b_2(S)=22$, and the rational
N\'eron--Severi space
\[
 N=\operatorname{NS}(S)\otimes\mathbb Q\subset H^2(S,\mathbb Q)
\]
has dimension $\rho\le20$. Its intersection form is nondegenerate.
Choose divisor classes $D_1,\ldots,D_\rho$ forming a basis of $N$, let
$G_{ij}=\int_S D_iD_j$, and write $(G^{ij})=G^{-1}$. Set
\[
 T=N^\perp,\qquad H^2(S,\mathbb Q)=N\oplus T,
 \qquad \dim T=22-\rho>0.
\]
Let $o$ denote the cohomology class of a closed point, normalized by
$\int_S o=1$, and use $a\boxtimes b=\operatorname{pr}_1^*a\smile
\operatorname{pr}_2^*b$. Consider the explicitly algebraic rational class
\begin{equation}\label{cont4:tau}
 \tau=[\Delta_S]-o\boxtimes1-1\boxtimes o
                 -\sum_{i,j}G^{ij}D_i\boxtimes D_j.
\end{equation}
Every term has codimension two on $X$. Products of divisor classes are
interpreted in the Chow ring, so no transverse representatives need be
assumed without moving them. The inverse matrix has rational entries,
which is exactly where rational coefficients are used.

\textbf{Proposition 4.1.} The class $\tau$ is nonzero and is not in the
degree-four part of the algebra generated by rational divisor classes of $X$.

\emph{Proof.} For a degree-four correspondence $Z$, use the action
\[
 Z_*(v)=(\operatorname{pr}_2)_*
                 \bigl(\operatorname{pr}_1^*v\smile Z\bigr),
 \qquad v\in H^2(S,\mathbb Q).
\]
The diagonal acts as the identity. The two point terms act as zero on
$H^2$, by degree. The last term acts by
\[
 v\longmapsto\sum_{i,j}G^{ij}\left(\int_S vD_i\right)D_j,
\]
the orthogonal projector onto $N$. Hence $\tau_*$ is zero on $N$ and the
identity on $T$. Since $T\ne0$, $\tau\ne0$.

K\"unneth and $H^1(S)=0$ give
$H^2(X,\mathbb Q)=H^2(S,\mathbb Q)\oplus H^2(S,\mathbb Q)$.
The rational $(1,1)$ part is $N\oplus N$. Consequently all divisor
classes on $X$ are sums of pullbacks of classes in $N$. Their degree-four
products lie in
\begin{equation}\label{cont4:divspan}
 \mathbb Q(o\boxtimes1)\;\oplus\;(N\otimes N)\;
                       \oplus\;\mathbb Q(1\boxtimes o).
\end{equation}
Each correspondence in this space annihilates $T$. The action of $\tau$
on $T$ therefore proves the assertion.\hfill$\square$

This is a failure of a proposed generating set, not a counterexample to
the Hodge conjecture: equation~\eqref{cont4:tau} already exhibits $\tau$
as a rational algebraic cycle. It also shows why having a basis of the
divisor classes on each factor is insufficient on the product. The
diagonal contains a component that remembers the transcendental
cohomology of the factors. Here ``transcendental'' describes $T$ on $S$,
not the algebraicity of the correspondence on $S\times S$.

\paragraph{An explicit intersection check.}
For an ample divisor $h$ on $S$, put
$H=h\boxtimes1+1\boxtimes h$, an ample class on $X$. The preceding
projector description gives
\begin{equation}\label{cont4:degree}
 \int_X\tau H^2=0,\qquad
 \int_X\tau^2=22-\rho>0.
\end{equation}
For the first equality, expand $H^2$ into the two point-factor terms and
$2h\boxtimes h$. The point-factor terms pair trivially with the
$H^2\otimes H^2$ component $\tau$, and the middle term pairs to zero
because $h\in N$ whereas $\tau$ acts through $T$.

For the second equality, choose any rational basis $t_a$ of $T$ with
intersection matrix $B$. The K\"unneth expression for this component is
$\tau=\sum_{a,b}(B^{-1})_{ab}t_a\boxtimes t_b$. All relevant degrees
are even, so there is no Koszul sign. Direct contraction gives
\[
 \int_X\tau^2
 =\sum_{a,b,c,d}(B^{-1})_{ab}(B^{-1})_{cd}B_{ac}B_{bd}
 =\operatorname{tr}(I_T)=\dim T.
\]
This also independently detects nonvanishing without appealing only to
the correspondence action. It uses no positivity of the form on $T$.

\paragraph{Stress test: effective cycles are the wrong target.}
Any nonzero effective rational codimension-two cycle
$Z=\sum a_jZ_j$ with $a_j>0$ satisfies
\[
 \int_X[Z]H^2=\sum_j a_j\deg_H(Z_j)>0.
\]
Each degree is positive because $H$ is ample. Equation~\eqref{cont4:degree}
therefore proves that no positive rational multiple of $\tau$ can be
represented by a nonzero effective rational cycle. The zero cycle cannot
represent it either. Thus an algorithm that tries to represent every
rational Hodge class by one effective cycle, even after clearing
denominators, must fail on this example.

Allowing differences restores the correct formulation: the positive
and negative terms in \eqref{cont4:tau} can be collected separately.
Adding a large multiple of $H^2$ gives positive degree but does not, just
from that numerical fact, give an effective representative. Degree is
one linear functional on a space of classes; it is not a membership test
for the effective cone. An argument using this modification would have
to establish an actual effectivity theorem and a way to subtract the
added algebraic class afterwards.

\paragraph{A conditional continuation with a checkable endpoint.}
Let $A^p(X)\subset H^{2p}(X,\mathbb Q)$ be the span of algebraic classes.
The elementary strategy above proves algebraicity of the pieces obtained
from known lower-degree algebraic classes by multiplication by $h$.
It leaves the rational primitive Hodge pieces as the first uncontrolled
terms. A sufficient continuation is the following precise input: for
each such primitive piece, produce algebraic classes whose projections
span it, with the projections and all lower Lefschetz pieces already
known to be algebraic. Under these assumptions, subtract the known
lower pieces, span the primitive one, and add the pieces back. Induction
through the Lefschetz decomposition then gives the desired class.

The stated input is not proved here. In particular, it cannot be supplied
by assuming that all rational Hodge classes are products of divisors;
Proposition 4.1 disproves that. Nor does the existence of a harmonic
representative supply algebraic support or the required cycle. These
are geometrical existence questions beyond decomposition in cohomology.

\paragraph{Verification and outcome.}
The argument is symbolic: the matrix projector, its square, its action
on $T$, and both intersection numbers have been derived explicitly.
Only the listed standard Hodge and K3 facts are external inputs. There
is no assertion of a search over varieties, a numerical period
certificate, or a formal proof check. The original rational statement
is preserved throughout.

\textbf{Outcome: Partial argument and two disproved shortcuts; UNRESOLVED.}
The divisor propagation covers the standard low-dimensional cases.
On a K3 square the class $\tau$ gives an exact, nonzero algebraic
counterexample both to divisor generation and to universal effectivity.
What remains for the general problem is the construction of the missing
primitive algebraic classes, not a choice of coefficients or a better
intersection-matrix calculation. No resolution of the general Hodge
conjecture is claimed.

\subsection{Sources}
\begin{itemize}
\item[S1]Clay Mathematics Institute, The Millennium Prize Problems.
\url{https://www.claymath.org/library/monographs/MPPc.pdf}.
\textit{Formulation source.}
\item[S2]Hodge Conjecture.
\url{https://www.claymath.org/millennium/hodge-conjecture/}.
\textit{Further reference; not a proof endorsement.}
\item[S3]A Proof of the Hodge Conjecture Derived from One Semiregular Witness Axiom — Mohammad F Islam.
\url{https://zenodo.org/records/22701510}.
\textit{Further reference; not a proof endorsement.}
\item[S4]The Hodge conjecture for Fermat fourfolds of odd degree at most 199 — Rifat Jumagulov.
\url{https://arxiv.org/abs/2608.18134}.
\textit{Further reference; not a proof endorsement.}
\item[E1]Pierre Deligne, The Hodge Conjecture, Clay Mathematics Institute, Sections 1--2.
\url{https://www.claymath.org/wp-content/uploads/2022/06/hodge.pdf}.
\textit{Added primary source: Definitions and formulation. Consulted 24 September 2026.}
\item[E2]Daniel Huybrechts, \emph{Lectures on K3 Surfaces}, author manuscript, Chapters 1 and 3 (intersection form, N\'eron--Severi and transcendental spaces).
\url{https://www.math.uni-bonn.de/people/huybrech/K3Global.pdf}.
\textit{Primary author text consulted 27 September 2026. The projector and shortcut counterexamples are computed in this attempt.}
\end{itemize}
\end{document}
